\documentclass[12pt,a4paper]{article}
\usepackage[margin=25mm, headheight=15pt, headsep=10pt]{geometry}
\usepackage{fontspec}
\usepackage[table]{xcolor} % Note: table option is required for rowcolors
\usepackage{titlesec}
\usepackage{tocloft}
\usepackage{fancyhdr}
\usepackage{booktabs}
\usepackage{tcolorbox}
\tcbuselibrary{skins, breakable}
\usepackage[colorlinks=true, linkcolor=emerald, urlcolor=emerald, bookmarks=true]{hyperref}

\definecolor{emerald}{RGB}{0,102,51}
\definecolor{darkred}{RGB}{139,0,0}
\definecolor{lightgray}{RGB}{245,245,245}

\usepackage[nil]{babel}
\babelprovide[import, main]{bengali}
\babelprovide[import]{arabic}
\babelfont{rm}[Renderer=HarfBuzz,Scale=1.0]{Noto Serif Bengali}
\babelfont{sf}[Renderer=HarfBuzz]{Noto Sans Bengali}
\babelfont{tt}[Renderer=HarfBuzz]{Noto Sans Bengali}
\babelfont[arabic]{rm}[Renderer=HarfBuzz,Scale=1.1]{Noto Naskh Arabic}

% Section styling
\titleformat{\section}{\Large\bfseries\color{emerald}}{\thesection.}{0.5em}{}
\titleformat{\subsection}{\large\bfseries\color{emerald!85!black}}{\thesubsection.}{0.5em}{}
\titleformat{\subsubsection}{\normalsize\bfseries\color{emerald!70!black}}{\thesubsubsection.}{0.5em}{}
\titlespacing*{\section}{0pt}{18pt}{8pt}

% Fancy Header/Footer setup
\pagestyle{fancy}
\fancyhf{} % clear all header and footer fields
\fancyhead[L]{\color{emerald}\small\textbf{\nouppercase{\leftmark}}}
\fancyfoot[L]{\scriptsize\color{black!50}{\lat{AI}}-সংকলিত, আলেম-যাচাইকৃত নয়}
\fancyfoot[C]{\color{emerald!70!black}\thepage}
\renewcommand{\headrulewidth}{0.4pt}
\renewcommand{\headrule}{\hbox to\headwidth{\color{emerald}\leaders\hrule height \headrulewidth\hfill}}
\renewcommand{\footrulewidth}{0pt}

% Custom boxes
% 1. Ayat/Hadith Box
\newtcolorbox{ayatbox}[1][]{
    enhanced, breakable, 
    colback=emerald!5, 
    colframe=emerald, 
    boxrule=0.5pt, 
    arc=3pt, 
    left=10pt, right=10pt, top=10pt, bottom=10pt,
    #1
}

% 2. Quote/Translation Box
\newtcolorbox{quotebox}[1][]{
    enhanced, breakable,
    frame hidden,
    colback=lightgray,
    borderline west={3pt}{0pt}{emerald},
    left=15pt, right=10pt, top=10pt, bottom=10pt,
    sharp corners,
    #1
}

% 3. AI-generated notice box
\newtcolorbox{warnbox}[1][]{
    enhanced, breakable,
    colback=red!4,
    colframe=darkred,
    boxrule=0.8pt,
    arc=3pt,
    left=10pt, right=10pt, top=10pt, bottom=10pt,
    #1
}

% Commands
\newcommand{\arab}[1]{\foreignlanguage{arabic}{#1}}
\newcommand{\kib}[1]{\textcolor{emerald}{\textbf{#1}}}

% Latin (English) fallback: Noto Serif Bengali has NO Latin glyphs
\newfontfamily\latfont[Renderer=HarfBuzz]{Noto Serif}
\newcommand{\lat}[1]{{\latfont #1}}

\setlength{\parskip}{5pt}
\setlength{\parindent}{0pt}

% Fix for missing bullet in Noto Serif Bengali
\renewcommand{\labelitemi}{$\bullet$}



\begin{document}
\begin{center}{\Large\bfseries\color{emerald} ব্যবহারিক গাইড — ছবি, ফটো ও AI-ছবি নিয়ে সিদ্ধান্তের পথ}\end{center}
\vspace{1em}
\section*{ব্যবহারিক গাইড — ছবি, ফটো ও \lat{AI}-ছবি নিয়ে সিদ্ধান্তের পথ}

\begin{warnbox}
 \textbf{গুরুত্বপূর্ণ নোটিশ (\lat{Important} \lat{Notice}):} এই কন্টেন্টটি \lat{AI} (কৃত্রিম বুদ্ধিমত্তা) দিয়ে প্রস্তুত ও সংকলিত। \lat{AI} কঠোর সূত্র-মান নীতি মেনে শুধু নির্ভরযোগ্য উৎস থেকে তথ্য সংগ্রহ করেছে — কেবল সহীহ/হাসান হাদিস ও সহীহ সনদের আসার রাখা হয়েছে, দুর্বল সনদ বাদ দেওয়া হয়েছে। তবে এটি কোনো আলেম বা মুহাদ্দিস কর্তৃক যাচাইকৃত নয়।
\end{warnbox}

\begin{quote}
এই ফাইলটি দ্বিধার \textbf{\lat{practice} (ব্যবহারিক)} সমাধান দেখায় — শুধু তত্ত্ব নয়। প্রতিটি প্রশ্নের পাশে দলিল-নির্ভর উত্তর এবং আলেমদের মতভেদ সৎভাবে দেওয়া হলো।
\end{quote}

\medskip\hrule\medskip

\subsection*{১. দ্রুত উত্তর (\lat{Quick} \lat{Answers})}

\begin{table}[h]\centering\small
\begin{tabular}{ll}
প্রশ্ন & উত্তর \\
\hline
শুধু মানুষের ছবি আঁকা হারাম? & \textbf{না} — মানুষ, পশু, পাখিসহ \textbf{যেকোনো প্রাণবান (জীবনবিশিষ্ট) সৃষ্টির ছবি আঁকা/গড়া হারাম} (সহীহ বুখারি-মুসলিমের হাদিস) \\
পশুপাখির ছবি আঁকা জায়েজ? & \textbf{না} — প্রাণ আছে এমন সব প্রাণীর ছবিই নিষিদ্ধ; হাদিসে \lat{“}গাছের ছবি বানাও\lat{”} (মুসলিম ২১১০খ) \\
গাছ, পাহাড়, বাড়ি, নৌকার ছবি? & \textbf{জায়েজ} — জড়বস্তুর ছবিতে কোনো নিষেধ নেই (মুসলিম ২১১০খ) \\
ফটোগ্রাফি (মোবাইল-ক্যামেরার ছবি)? & \textbf{মতভেদ} — অধিকাংশ সমসাময়িক আলেম জায়েজ বলেছেন (প্রতিফলন ধারণ, তাসবীর নয়); কিছু আলেম হারাম বলেছেন \\
\lat{AI} দিয়ে ছবি বানানো? & \textbf{নতুন মতভেদ} — এক দল আঁকার মতো হারাম; অন্য দল ডিজিটাল টুল, শর্তসাপেক্ষে জায়েজ; সর্বসম্মত — জড়বস্তু ও মুখবিহীন ছবি জায়েজ \\
ঘরে ফ্রেম করে প্রাণীর ছবি টাঙানো? & \textbf{নিষিদ্ধ} — ফেরেশতা ঢোকে না (বুখারি ৫৯৫৭, মুসলিম ২১০৭); বালিশ-কুশন-কার্পেটে ছবি জায়েজ (মুসলিম ২১০৮) \\
\hline
\end{tabular}
\end{table}

\medskip\hrule\medskip

\subsection*{২. মূল নীতিটা এক লাইনে}

\begin{quote}
\textbf{ইসলাম ছবি-আঁকাকে নিষিদ্ধ করে না — নিষিদ্ধ করে প্রাণবান সৃষ্টির \lat{“}নকল\lat{”} (\lat{imitation}) করাকে।}
\end{quote}

যেকোনো সিদ্ধান্তে এই নীতিটি ফিরে আসুন:
\begin{itemize}
\item \textbf{প্রাণ আছে (আঁকা/গড়া)} $\rightarrow$ হারাম
\item \textbf{প্রাণ নেই (আঁকা/গড়া)} $\rightarrow$ জায়েজ
\item \textbf{প্রাণ আছে, কিন্তু ফটো (বাস্তবের ধারণ)} $\rightarrow$ মতভেদ, অধিকাংশে জায়েজ
\item \textbf{প্রাণ আছে, কিন্তু \lat{AI}-তৈরি} $\rightarrow$ মতভেদ (নিচে বিস্তারিত)
\item \textbf{প্রাণ আছে, অবজ্ঞার জিনিসে (বালিশ-কার্পেট)} $\rightarrow$ জায়েজ (অধিকাংশ মত)
\item \textbf{প্রাণ আছে, ঘরে ঝুলিয়ে সম্মান} $\rightarrow$ হারাম/সতর্কতা
\end{itemize}
\medskip\hrule\medskip

\subsection*{৩. মানুষ বনাম পশুপাখি — দ্বিধার সরাসরি উত্তর}

\textbf{দ্বিধা:} \lat{“}ইসলামে ছবি আঁকা জায়েজ নেই\lat{”} — এখানে কি শুধু মানুষের ছবি, নাকি পশুপাখিরও?

\textbf{উত্তর:} সহীহ হাদিসে নিষেধ মানুষের মধ্যে \textbf{সীমাবদ্ধ নয়}। হাদিসগুলোতে শব্দ হলো \lat{“}আল-মুসাওয়িরুন\lat{”} (ছবি নির্মাতারা) — সাধারণ শব্দ। ইবনে আব্বাস (রাঃ) শিল্পীকে বলেছেন: \lat{“}গাছের ছবি বানাও, এমন বস্তুর ছবি বানাও যাতে প্রাণ নেই\lat{”} (মুসলিম ২১১০খ)। অর্থাৎ \textbf{যেকোনো প্রাণবান সৃষ্টি} — মানুষ হোক, ঘোড়া, বাঘ, পাখি, মাছ — সবই নিষেধের অন্তর্ভুক্ত।

\textbf{কেন?} কারণ নিষেধের যুক্তি হলো প্রাণ সৃষ্টির দাবি। মানুষ আঁকা-গড়ায় যতই দক্ষ হোক, সে প্রাণ দিতে পারে না। এই যুক্তি মানুষের ছবিতে যেমন, পশুপাখির ছবিতেও তেমনই।

\textbf{সতর্কতা:} পশুপাখির ছবি আঁকা \lat{“}জায়েজ\lat{”} ভাবার কোনো সুযোগ নেই — এটি দুর্বল বা অদ্ভুত মত; প্রায় সব আলেম এতে একমত। (হানাফি মাযহাবে মাথাবিহীন/অসম্পূর্ণ ছবির ক্ষেত্রে শিথিলতা আছে — সেটি আলাদা আলোচনা, নিচের টেবিলে।)

\medskip\hrule\medskip

\subsection*{৪. ফটোগ্রাফি — আঁকার মতো নয় (মতভেদ)}

ফটোগ্রাফির বিধান নিয়ে আলেমদের মধ্যে মতভেদ আছে। এটা মনে রাখা জরুরি — \textbf{এটি খাঁটি মাসআলার ইখতিলাফ}, দুই পক্ষেরই দলিল আছে।

\subsubsection*{মত ১ — ফটোগ্রাফি জায়েজ (অধিকাংশ সমসাময়িক আলেম)}

\textbf{যুক্তি:} ফটো ছবি \textbf{\lat{“}আঁকে\lat{”}} না — বাস্তব বস্তুর আলো-প্রতিফলনকে \textbf{ধারণ (\lat{capture})} করে। তাই এটি তাসবীর (নকল-সৃষ্টি) নয়, বরং দর্পণের প্রতিফলনের মতো।

\begin{itemize}
\item শাইখ ইবনে উসাইমীন (রহ.): ফটোগ্রাফি তাসবীরের অন্তর্ভুক্ত নয় — এটি বাস্তবকে ধারণ; জায়েজ।
\item শাইখ আল-আলবানী (রহ.): একই মত।
\item আযহার/দারুল ইফতা, জর্ডান, ইসলামওয়েব: ডিজিটাল ছবি জায়েজ।
\end{itemize}
\textbf{শর্ত:} ফটোতে হারাম কনটেন্ট (আওরাহ প্রকাশ, পাপের দৃশ্য) না থাকা; ঝুলিয়ে-টাঙিয়ে সম্মান না করা।

\subsubsection*{মত ২ — ফটোগ্রাফি হারাম (কঠোর মত)}

\textbf{যুক্তি:} ফটোও প্রাণবান সৃষ্টির বাস্তব প্রতিকৃতি; নিষেধের মূল — প্রাণবান সৃষ্টির প্রতিকৃতি, যেকোনো মাধ্যমে।

\textbf{পরামর্শ:} আপনার যে আলেমকে বিশ্বাস করেন, তার মত অনুসরণ করুন। নিরাপদ পথ — \textbf{পার্থক্যের জায়গায় ঘরের দেয়ালে ছবি টাঙানো এড়ানো}, যেটা দুই মতেই সতর্কতা।

\medskip\hrule\medskip

\subsection*{৫. \lat{AI}-দিয়ে ছবি বানানো — নতুন মাসআলা (নতুন মতভেদ)}

\lat{AI} (কৃত্রিম বুদ্ধিমত্তা) দিয়ে ছবি বানানো একটি \textbf{নতুন (\lat{novel}) বিষয়} — নবীজির (সা.) যুগে এর অস্তিত্ব ছিল না। তাই এটা নিয়ে আলেমরা এখনো মত দিচ্ছেন; কোনো একক \lat{“}সবাই মানা\lat{”} মত এখনো নেই। নিচে বর্তমান ফতোয়া-জগতের তিনটি অবস্থান সৎভাবে দেওয়া হলো:

\subsubsection*{অবস্থান ১ — \lat{AI}-প্রাণীর ছবি হারাম (কঠোর মত)}

\textbf{মত:} \lat{AI} দিয়ে মানুষ-পশুর ছবি বানানো আঁকার মতোই — তাসবীরের অন্তর্ভুক্ত; \lat{intention} বা ব্যবহার যাই হোক, হারাম। কারণ \lat{AI} যতই \lat{“}বুদ্ধিমান\lat{”} হোক, এটি প্রাণবান সৃষ্টির \textbf{নতুন নকল} বানাচ্ছে।

\textbf{উৎস:} দারুল ইফতা নিউ ইয়র্ক (\lat{askthemufti}.\lat{us}) — \lat{AI}-তৈরি মানুষ/পশুর ছবি-ভিডিও তাসবীরের অন্তর্ভুক্ত বলে স্পষ্ট ফতোয়া দিয়েছে। সুবিধা — \textbf{\lat{faceless} (মুখবিহীন) ছবি জায়েজ}।

\subsubsection*{অবস্থান ২ — \lat{AI}-ছবি আঁকার মতো নয়, শর্তসাপেক্ষে জায়েজ (নমনীয় মত)}

\textbf{যুক্তি:} \lat{AI} একটি টুল/প্রোগ্রাম; ডিজিটাল ছবির শারীরিক অস্তিত্ব নেই (স্ক্রিনে দেখা যায়); হাতে আঁকার মতো \lat{“}তাসবীর\lat{”} নয়; মূল নিয়ন্ত্রণ বিষয়বস্তু ও ব্যবহারের।

\begin{itemize}
\item \textbf{জর্ডান দারুল ইফতা (২০২৫):} ফটোকে \lat{AI}-রূপান্তরের মত — তাসবীর নয়; তবে বিষয়বস্তু বৈধ হতে হবে (আওরাহ নয়, উপহাস নয়, প্রতারণা নয়)।
\item \textbf{ইসলামওয়েব (২০২৩):} \lat{AI}-তৈরি মানুষ-প্রাণীর ইলেকট্রনিক ছবি নিষিদ্ধ নয় — কারণ এর বাস্তব সত্তা নেই।
\item \textbf{দারুল ইফতা বার্মিংহাম:} \lat{AI} ডিজিটাল আঁকা-তুলনার থেকে আলাদা; শিক্ষা/প্রয়োজনে জায়েজ; তবে \textbf{বাস্তবধর্মী (\lat{realistic})} মানুষের ছবি বেশি সমস্যাযুক্ত — নমনীয় পদ্ধতি।
\item \textbf{সিকার্স গাইডেন্স (শাফেয়ি):} টুল নিজে হারাম নয় — যা বানাও তা-ই বিধান নির্ধারণ করে; জড়বস্তু সব মতেই জায়েজ।
\end{itemize}
\subsubsection*{অবস্থান ৩ — সব পক্ষের সম্মত শর্ত (\lat{Common} \lat{Ground})}

নিচের বিষয়গুলো \textbf{সব মতেই} নিষিদ্ধ — এগুলো নিয়ে মতভেদ নেই:

\begin{enumerate}
\item \textbf{আওরাহ/নগ্নতা/যৌন কনটেন্ট} — কোনো মতেই নয়।
\item \textbf{উপহাস-বিদ্রূপ (\lat{mockery})} — সূরা হুজুরাত ৪৯:১১-এর লঙ্ঘন।
\item \textbf{প্রতারণা-জালিয়াতি (\lat{deception}/\lat{fraud})} — মিথ্যা ছবি, ভুয়া পরিচয়।
\item \textbf{হারাম দৃশ্য প্রচার} — পাপের প্রচারে \lat{AI} ব্যবহার।
\item \textbf{জড়বস্তুর \lat{AI} ছবি} — সব মতেই জায়েজ (গাছ, প্রকৃতি, ভবন, অ্যাবস্ট্রাক্ট আর্ট)।
\end{enumerate}
\subsubsection*{আপনার সিদ্ধান্তের পথ (\lat{Decision} \lat{Path})}

একটি সহজ \textbf{\lat{flowchart} (প্রবাহচিত্র)} মাথায় রাখুন:

```
আপনি \lat{AI} দিয়ে ছবি বানাতে চান
├── জড়বস্তু/প্রকৃতি/অ্যাবস্ট্রাক্ট? $\rightarrow$ ✅ জায়েজ (সব মত)
├── মানুষ/প্রাণীর ছবি?
│   ├── মুখবিহীন (\lat{faceless})? $\rightarrow$ ✅ জায়েজ (কঠোর মতেও)
│   ├── মুখসহ?
│   │   ├── শিক্ষা/চিকিৎসা/প্রয়োজনীয়? $\rightarrow$  নমনীয় মতের ফতোয়া অনুসরণযোগ্য
│   │   └── শখ/সোশ্যাল মিডিয়া? $\rightarrow$  কঠোর মত মতে হারাম; তাকওয়া = বিরত থাকা
└── কনটেন্টে হারাম কিছু (আওরাহ/উপহাস/প্রতারণা)? $\rightarrow$ ❌ হারাম (সব মত)
```

\subsubsection*{তাকওয়ার সুপারিশ (\lat{Recommended} \lat{Path})}

সব মতের আলোকে সবচেয়ে \textbf{নিরাপদ (\lat{safe})} পথ:

\begin{enumerate}
\item \lat{AI} দিয়ে মানুষ-প্রাণীর \textbf{বাস্তবধর্মী মুখসহ ছবি} এড়িয়ে চলুন — এটিই একমাত্র বিষয় যেখানে মতভেদ আছে এবং কঠোর মত হারাম বলেছে।
\item দরকার হলে \textbf{\lat{faceless}}, কার্টুনিশ বা \textbf{জড়বস্তুর} ছবি ব্যবহার করুন — সব মতেই জায়েজ।
\item শিক্ষা, চিকিৎসা (মেডিকেল ইলাস্ট্রেশন), ডকুমেন্টারি-প্রয়োজনে যদি জরুরি হয়, তাহলে নমনীয় মত অনুসরণ করুন — তবে বিনা প্রয়োজনে নয়।
\item \lat{AI} দিয়ে কোনো মানুষকে বাস্তবে নেই এমন দৃশ্যে বসানো (গভফেক), কারো চেহারা অপব্যবহার — \textbf{নিষিদ্ধ} (প্রতারণা)।
\end{enumerate}
\medskip\hrule\medskip

\subsection*{৬. ঘরের সাজসজ্জা — কি রাখা, কি না}

\begin{table}[h]\centering\small
\begin{tabular}{lll}
বস্তু & বিধান & দলিল \\
\hline
দেয়ালে-দরজায় ফ্রেম করা প্রাণীর ছবি & \textbf{নিষিদ্ধ} — ফেরেশতা ঢোকে না & বুখারি ৫৯৫৭, মুসলিম ২১০৭ \\
আয়াত/নবীজির নামের ক্যালিগ্রাফি & \textbf{জায়েজ} — এটি ছবি নয়, লেখা & সাধারণ নীতি \\
বালিশ, কুশন, কার্পেট, গদিতে ছবি & \textbf{জায়েজ} (অবজ্ঞার বস্তু) & মুসলিম ২১০৮ \\
মোবাইল-কম্পিউটারের ফটো (ডিজিটাল) & মতভেদ — অধিকাংশে জায়েজ & সমসাময়িক ফতোয়া \\
মেয়েদের পুতুল-খেলনা (ছোট শিশু) & মতভেদ — বুখারি ৬১৩০-এর আলোকে অনেকে জায়েজ; বড় বাস্তবধর্মী মূর্তি সতর্কতা & বুখারি ৬১৩০ \\
টিভি-পর্দায় চলমান ছবি & মতভেদ — জায়েজ (ইবনে উসাইমীনসহ অনেকে) & সমসাময়িক ফতোয়া \\
\hline
\end{tabular}
\end{table}

\begin{quote}
\textbf{টিভি/স্ক্রিন নোট:} স্ক্রিনে ছবি দেখা নিয়ে অধিকাংশ সমসাময়িক আলেম জায়েজ বলেছেন — কারণ এটি ঝুলিয়ে রাখার মতো স্থায়ী তাসবীর নয়। তবে স্ক্রিনে আওরাহ/হারাম দৃশ্য আলাদা বিষয় — সেটি নিষিদ্ধ।
\end{quote}

\medskip\hrule\medskip

\subsection*{৭. খেলনা-পুতুল — শিশুদের জন্য}

আয়েশা (রাঃ) নবীজির (সা.) কাছে পুতুল নিয়ে খেলতেন — সহীহ বুখারি ৬১৩০। এর ভিত্তিতে অনেক আলেম ছোট মেয়েদের পুতুল-খেলনাকে অনুমোদন করেছেন। তবে:

\begin{itemize}
\item \textbf{মতভেদ আছে} — কেউ কেউ এটিকে শুধু সেই যুগ ও সেই বয়সের জন্য সীমিত মনে করেন।
\item \textbf{নিরাপদ পথ:} ছোট শিশুদের জন্য সাধারণ পুতুল/খেলনা রাখা যেতে পারে (হানাফি-হাম্বলির অনুমতি); তবে বড়, বাস্তবধর্মী, মূল্যবান বা পূজার মতো রাখা মূর্তি — এড়িয়ে চলুন।
\item শিশু খেলনার উদ্দেশ্য যখন বাস্তবধর্মী \textbf{\lat{showing} (প্রদর্শনী)} নয়, তখন আলেমদের প্রশস্ততা বেশি।
\end{itemize}
\medskip\hrule\medskip

\subsection*{৮. পুরনো ছবি ও তাওবার নিয়ম}

যদি আগে প্রাণীর ছবি আঁকা/টাঙানো হয়ে থাকে:

\begin{enumerate}
\item \textbf{টাঙানো ছবি:} নামিয়ে দিন, ধ্বংস করুন বা ব্যবহার্য (বালিশ-কাপড়) বস্তুতে রূপান্তর করুন — নবীজির (সা.) আমলের মতো।
\item \textbf{আঁকা ছবি:} মুছে ফেলুন, বিশেষত \textbf{মাথা/চেহারা} মুছে ফেললে অনেক আলেমের মতে ছবি নিষিদ্ধ অবস্থা থেকে বেরিয়ে যায় (হানাফি পদ্ধতি)।
\item \textbf{ডিজিটাল ফটো:} মুছে ফেলতে বাধ্য নন (অধিকাংশ মত) — তবে হারাম কনটেন্ট মুছুন; জায়েজ বিষয়ের ছবি রাখা যেতে পারে।
\item \textbf{তাওবা:} আন্তরিক ইস্তিগফার করুন — অতীতের জন্য আল্লাহর কাছে ক্ষমা; সামনে থেকে বিরত থাকা। আল্লাহ ক্ষমাশীল।
\end{enumerate}
\medskip\hrule\medskip

\subsection*{৯. সিদ্ধান্তের সারসংক্ষেপ — মনে রাখার কার্ড}

\begin{quote}
\textbf{১. প্রাণবান সৃষ্টির ছবি আঁকা/গড়া = হারাম} (মানুষ + পশু + পাখি — সব) \\
\textbf{২. জড়বস্তুর ছবি = জায়েজ} (গাছ, পাহাড়, ভবন, নৌকা) \\
\textbf{৩. ঘরে ঝুলিয়ে-টাঙিয়ে প্রাণীর ছবি = নিষিদ্ধ} (ফেরেশতা ঢোকে না) \\
\textbf{৪. বালিশ-কুশন-কার্পেটে ছবি = জায়েজ} \\
\textbf{৫. ফটোগ্রাফি = মতভেদ; অধিকাংশে জায়েজ} — ঝুলিয়ে সম্মান না করলেই নিরাপদ \\
\textbf{৬. \lat{AI}-ছবি = নতুন মতভেদ; মুখসহ বাস্তবধর্মী প্রাণীর ছবি এড়ানো = তাকওয়া; \lat{faceless}/জড়বস্তু = জায়েজ} \\
\textbf{৭. যেকোনো মাধ্যমেই আওরাহ, উপহাস, প্রতারণা = হারাম}
\end{quote}

\medskip\hrule\medskip

\subsection*{১০. রেফারেন্স}

\begin{itemize}
\item সহীহ বুখারি: ৫৯৫০, ৫৯৫১, ৫৯৫৪, ৫৯৫৭, ৫৯৬১, ৫৯৬৩, ৩২২৬, ৬১৩০
\item সহীহ মুসলিম: ২১০৬, ২১০৭, ২১০৮, ২১০৯, ২১১০খ
\item সমসাময়িক ফতোয়া: জর্ডান দারুল ইফতা (আলিফতাআ), ইসলামওয়েব, দারুল ইফতা বার্মিংহাম, দারুল ইফতা নিউ ইয়র্ক, সিকার্স গাইডেন্স
\item ইবনে উসাইমীন ও আল-আলবানী (রহ.) — ফটোগ্রাফি সংক্রান্ত মত
\item এই ফাইলটি \lat{AI}-সংকলিত, আলেম-যাচাইকৃত নয় — গুরুত্বপূর্ণ সিদ্ধান্তের আগে একজন বিশ্বস্ত আলেমের সাথে \lat{consult} (পরামর্শ) করুন।
\end{itemize}

\end{document}
